% Unit 5 map -- one page.
\documentclass{apchem-worksheet}

\apcunit{5}
\apcunittitle{Kinetics}
\apcdoctitle{Unit Map}
\apcsubtitle{CED Unit 5 \textbullet\ 7--9\% of the AP Exam \textbullet\ 7 blocks + float + exam}

\begin{document}
\apctitleblock

\begin{apcnote}[How this unit is graded --- and what makes it different]
Unit exam \textbf{Fri Nov 20}: 25 multiple choice (\textbf{no calculator})
$+$ 1 short and 1 long free-response, 69 minutes, 39 points.

Every other unit so far has asked \emph{what} happens. This one asks
\emph{how fast}, and the answer can never be read off a balanced equation.
\textbf{Orders come from experiment, always} --- the single exception being
an elementary step, where molecularity does give the order. Knowing which
case you are in is most of the unit.
\end{apcnote}

\section*{\sffamily Block calendar}

\renewcommand{\arraystretch}{1.3}
\noindent\begin{tabular}{@{}p{2.1cm}p{5.0cm}p{1.5cm}p{3.9cm}p{1.9cm}@{}}
\toprule
\textbf{Date} & \textbf{Topics} & \textbf{CED} & \textbf{Read (PDF pp.)} & \textbf{Practice}\\
\midrule
Mon Nov 2 & Reaction rates; measuring them & 5.1 & \S12.1 \textbullet\ 594--598 & WS 1\\
Wed Nov 4 & Rate laws and orders from data & 5.2 & \S12.2--12.3 \textbullet\ 598--604 & WS 1\\
Fri Nov 6 & Concentration over time; half-life & 5.3 & \S12.4 \textbullet\ 604--615 & WS 2\\
Mon Nov 9 & Elementary steps and mechanisms & 5.4, 5.7, 5.8 & \S12.5 \textbullet\ 615--620 & WS 3\\
Wed Nov 11 & Collision model; energy profiles & 5.5, 5.6 & \S12.6 \textbullet\ 620--626 & WS 3\\
Fri Nov 13 & Pre-equilibrium; multistep profiles & 5.9, 5.10 & \S12.5--12.6 \textbullet\ 615--626 & WS 4\\
Mon Nov 16 & Catalysis & 5.11 & \S12.7 \textbullet\ 626--632 & WS 4 (FRQ)\\
Wed Nov 18 & Reteach / float & --- & review sheet & ---\\
\textbf{Fri Nov 20} & \textbf{UNIT 5 EXAM} & all & review sheet & ---\\
\bottomrule
\end{tabular}

{\sffamily\footnotesize Dates from the co-op calendar (Thanksgiving week
Nov 23 is off). Shift if the schedule slips.}

\vspace{0.5em}
\section*{\sffamily By the end of this unit you can\ldots}

\begin{itemize}[leftmargin=*,itemsep=0.12em]
  \item Express a reaction rate in terms of any reactant or product, using
        the coefficients correctly. \ced{5.1}
  \item Determine a rate law and rate constant from initial-rate data.
        \ced{5.2}
  \item Identify reaction order from which graph is linear, and use the
        integrated rate law and half-life. \ced{5.3}
  \item Write the rate law for an elementary step from its molecularity.
        \ced{5.4}
  \item Explain rate in terms of collision frequency, energy, and
        orientation. \ced{5.5}
  \item Read and draw a reaction energy profile, identifying $E_a$ and
        $\Delta H$. \ced{5.6}
  \item Evaluate whether a proposed mechanism is consistent with the
        overall equation and the observed rate law. \ced{5.7}, \ced{5.8}
  \item Derive a rate law when the first step is a fast equilibrium.
        \ced{5.9}
  \item Interpret a multistep energy profile and identify the
        rate-determining step. \ced{5.10}
  \item Explain how a catalyst changes the rate without changing $\Delta H$
        or $K$. \ced{5.11}
\end{itemize}

\begin{apctrap}[Where students lose points in this unit]
Reading orders off the balanced equation; forgetting the coefficient when
converting between reactant and product rates; calling a graph ``linear''
without saying \emph{which} quantity was plotted; using $t_{1/2} = 0.693/k$
for a reaction that is not first order; leaving an intermediate in a final
rate law; and saying a catalyst ``shifts the equilibrium'' or ``makes the
reaction more favorable.''
\end{apctrap}

\end{document}
